% =================================================================
\chapter{Migration}
% =================================================================
\label{chap:migration}

In diesem Kapitel wird beschrieben, wie die neue Logging-Strategie bei Kunden umgesetzt werden kann.  
Die Migration erfolgt in mehreren Schritten, um Stabilität, Konsistenz und Nachvollziehbarkeit sicherzustellen.

\section{Betroffene Server und Services}
\begin{itemize}
    \item Alle betroffenen Server und Services identifizieren.
    \item Prüfen, welche Log4j-Versionen pro Service im Einsatz sind.
\end{itemize}

\section{Prüfung der Log-Ordner}
\begin{itemize}
    \item Existenz der Log-Ordner prüfen.
    \item Schreibrechte für den Service-User sicherstellen.
    \item Verfügbaren Speicherplatz prüfen.
\end{itemize}

\section{Rollover-Ordner}
\begin{itemize}
    \item Existenz des Rollover-Ordners prüfen.
    \item Festlegen, welche Komprimierung für Rollover-Dateien genutzt wird.
    \item (Optional) Netzwerkverfügbarkeit von Archivservern prüfen (z.B. Mounts, NFS, Freigaben).
\end{itemize}

\section{Pfadkonzept}
\begin{itemize}
    \item Entscheidung treffen: Symlinks oder globale Pfade.
    \item Globale Pfade vereinfachen das Setup, da lokale Ordner entlastet werden. Die Konfiguration wird dadurch nur marginal komplexer.
    \item Bei Symlinks: korrekte Auflösung prüfen, um Probleme zu vermeiden, wenn in applikationsbezogene Ordner geschrieben wird.
\end{itemize}

\section{Anpassung der Log-Konfigurationen}
Templates sind in Abschnitt \textit{\nameref{chap:logs_and_attachements}} hinterlegt.
\begin{itemize}
    \item Cron-Jobs oder andere Archivierungsjobs prüfen und berücksichtigen.
    \item Templates an lokale bzw. kundenspezifische Besonderheiten anpassen.
    \item Log-Pfade und Rollover-Pfade konfigurieren.
    \item Rollover-Policies prüfen (größe-basiert oder zeitbasiert).
    \item Rollover-Kompression setzen.
    \item Löschfristen für Logs definieren.
    \item Cron-Job für das Löschen leerer Datumsordner einrichten (siehe Abschnitt \ref{sec:cleanup}).
\end{itemize}

\section{Migration (pro Service)}
\begin{enumerate}
    \item Bereitstellen der neuen Konfigurationsdateien.
    \begin{enumerate}
        \item Pro Service eine separate Datei verwenden.
        \item Lokale Sicherung der alten Konfigurationen für schnelle Rollbacks erstellen.
        \item Alte Konfigurationsdatei durch die neue ersetzen.
        \item Prüfen, ob der Service bereits auf die neue Log-Struktur schreibt.
    \end{enumerate}
\end{enumerate}

\section{Verifikation}
\begin{itemize}
    \item Laufende Logs prüfen (Pfad, Format).
    \item Rollover überprüfen (Pfad, Kompression, täglicher Rollover).
\end{itemize}

\section{Cleanup}
\label{sec:cleanup}
\begin{itemize}
    \item Alte Logs in einen \texttt{\_Old}-Ordner im Rollover-Verzeichnis der jeweiligen Applikation verschieben.
    \begin{itemize}
        \item Beispiel: \texttt{mLogistics server.log.gz} $\rightarrow$ \texttt{/archiv/mLogistics/server/\_Old}
    \end{itemize}
    \item Alte Ordnerstrukturen entfernen.
    \item Nach 1–2 Rollovern die alten Strukturen endgültig prüfen und ggf. Konfigurationen auf Fehler kontrollieren.
\end{itemize}

\subsection{Cron-Job: Leere Datumsordner löschen}
Nach dem täglichen Rollover werden die Logs in den Datumsordner (\texttt{yyyy-mm-dd}) gelöscht wenn die jeweilige Frist überschritten ist.
Log4j sieht keine löschen dieser ordner vor weshalb ein Script genutzt werden muss.

Hierfür steht das Script \texttt{deleteEmptyFolders.sh} bereit, das rekursiv alle leeren Ordner
mit dem Namensmuster \texttt{yyyy-mm-dd} unterhalb eines angegebenen Verzeichnisses sucht und löscht.

\paragraph{Verwendung}
\begin{verbatim}
# Nur auflisten (kein Löschen)
./deleteEmptyFolders.sh /els/archiv/

# Auflisten und löschen
./deleteEmptyFolders.sh /els/archiv/ --delete
\end{verbatim}

\paragraph{Cron-Job einrichten}
Der Cron-Job wird mit \texttt{crontab -e} eingerichtet und läuft täglich um 02:00 Uhr. Je nach Systemumgebung muss statt mlogte01 ein anderer User genutzt werden. Es ist zu empfehlen die erste Node in einem Cluster zu nutzen und von dort die leeren Ordner aller nodes zu löschen:
\begin{verbatim}
0 2 * * * /els/mlogte01/deleteEmptyFolders.sh /els/archiv/ --delete >> /els/mlogte01/local/log/deleteEmptyFolders.log 2>&1
\end{verbatim}

Die Ausgabe wird in \texttt{deleteEmptyFolders.log} protokolliert. Mit \texttt{2>\&1} werden
auch Fehlermeldungen in dieselbe Logdatei umgeleitet.

\section{Rollback-Plan}
\begin{itemize}
    \item Schritte dokumentieren, um alte Konfigurationen und Logs wiederherzustellen.
    \item Alte Konfigurationen lokal und auf dem Server sichern.
    \item Bei Problemen: Alte Konfiguration wieder einspielen und ggf. Service neu starten.
    \item Alte Logs zurückverschieben, falls notwendig.
\end{itemize}